\documentclass[12pt,a4paper]{article}
\usepackage[a4paper,margin=18mm]{geometry}
\usepackage[T2A]{fontenc}
\usepackage[utf8]{inputenc}
\usepackage[russian]{babel}
\usepackage{amsmath,amssymb,mathtools}
\usepackage{array,booktabs,longtable}
\usepackage{xcolor}
\usepackage{hyperref}
\usepackage{tikz}

\hypersetup{
  colorlinks=true,
  linkcolor=blue!55!black,
  urlcolor=blue!55!black
}
\setlength{\parindent}{0pt}
\setlength{\parskip}{0.45em}
\setlength{\tabcolsep}{3pt}
\renewcommand{\arraystretch}{1.18}
\newcolumntype{P}[1]{>{\raggedright\arraybackslash}p{#1}}

\begin{document}
\hypersetup{pageanchor=false}

\begin{titlepage}
\thispagestyle{empty}
\begin{center}
\vspace*{0.22\textheight}
{\LARGE\bfseries Ревью домашней работы\par}
\vspace{1.4em}
{\large Автор: Лёвин Артём Александрович\par}
\vspace{0.8em}
{\large 12 сентября 2026 года\par}
\vfill
\begin{tikzpicture}
  \draw[line width=0.7pt]
    (0,0) .. controls (1.8,0.55) and (4.2,-0.55) .. (6,0);
\end{tikzpicture}
\end{center}
\end{titlepage}
\hypersetup{pageanchor=true}

\newpage
\section*{Сводка по заданиям с ошибками}
\small
\begin{center}
\begin{tabular}{P{0.07\linewidth} P{0.20\linewidth} P{0.27\linewidth} P{0.15\linewidth} P{0.15\linewidth}}
\toprule
№ задания & Тема & Суть ошибки & Ваш ответ & Правильный ответ \\
\midrule
\hyperref[task:4]{4} & Среднее арифметическое двух чисел & При сложении $x+3x$ получено $2x\cdot3$, то есть $6x$, вместо $4x$. & $12$ и $36$ & $18$ и $54$ \\
\addlinespace
\hyperref[task:8]{8} & Три последовательных натуральных числа & Условие о сумме заменено делением $3a$ на $3$; не учтены числа $a+1$ и $a+2$. & $114,115,116$ & $37,38,39$ \\
\bottomrule
\end{tabular}
\end{center}
\normalsize

\newpage
\thispagestyle{empty}
\vspace*{\fill}
\begin{center}
{\Huge\bfseries Разбор ошибок}
\end{center}
\vspace*{\fill}
\newpage

\phantomsection\label{task:4}
\section*{Задание 4}

\textbf{Текст задания.}

Среднее арифметическое двух положительных чисел равно $36$. Одно число в $3$ раза больше другого. Найдите оба числа.

\textbf{Ваше решение.}

В записи правильно выбрано неизвестное $x$ и второе число выражено как $3x$. Далее записано условие среднего арифметического:
\[
\frac{x+3x}{2}=36.
\]
После этого в работе появляется переход
\[
2x\cdot3=36\cdot2,
\]
а затем:
\[
2x\cdot3=72,
\qquad
2x=72:3,
\qquad
2x=24,
\qquad
x=12.
\]
Второе число найдено как
\[
12\cdot3=36.
\]

\textbf{Ваш ответ.}
\[
12 \text{ и } 36.
\]

\textcolor{red}{Ошибка:} неверно сложены подобные слагаемые $x$ и $3x$.

\textbf{Где именно возникает ошибка.}

Последний правильный шаг:
\[
\frac{x+3x}{2}=36.
\]

Первый неверный шаг:
\[
2x\cdot3=36\cdot2.
\]

Проблема находится в левой части. Сумма
\[
x+3x
\]
не равна
\[
2x\cdot3.
\]

\textbf{Почему этот шаг неверен.}

Слагаемые $x$ и $3x$ являются подобными: у них одинаковая буквенная часть $x$. При сложении подобных слагаемых складываются только коэффициенты при $x$:
\[
x+3x=1x+3x=(1+3)x=4x.
\]

Запись $2x\cdot3$ означает другое действие:
\[
2x\cdot3=6x.
\]
То есть вместо суммы $x+3x=4x$ в решение фактически попало выражение $6x$. Из-за этого значение $x$ получилось слишком маленьким.

Нарушено правило сложения подобных слагаемых: коэффициенты нужно \emph{складывать}, а не сначала заменять два слагаемых выражением $2x$, а затем умножать его на коэффициент $3$.

\textcolor{blue!60!black}{Ключевая идея:} если меньшее число обозначено через $x$, то большее равно $3x$, а их сумма равна $x+3x=4x$. Среднее арифметическое получается делением этой суммы на $2$.

\textbf{Исправление от последнего правильного шага.}

Возвращаемся к верному уравнению:
\[
\frac{x+3x}{2}=36.
\]
Сначала складываем подобные слагаемые:
\[
\frac{4x}{2}=36.
\]
Делим $4x$ на $2$:
\[
2x=36.
\]
Отсюда
\[
x=18.
\]
Тогда второе число:
\[
3x=3\cdot18=54.
\]

\textbf{Полное правильное решение.}

Пусть меньшее из двух положительных чисел равно $x$. По условию второе число в $3$ раза больше, значит оно равно $3x$.

Среднее арифметическое двух чисел --- это их сумма, делённая на $2$. Поэтому:
\[
\frac{x+3x}{2}=36.
\]
Складываем подобные слагаемые:
\[
\frac{4x}{2}=36.
\]
Получаем:
\[
2x=36.
\]
Делим обе части уравнения на $2$:
\[
x=18.
\]
Следовательно, меньшее число равно $18$, а большее:
\[
3\cdot18=54.
\]

\textbf{Проверка.}

Во-первых, $54$ действительно в $3$ раза больше $18$:
\[
54:18=3.
\]
Во-вторых, проверяем среднее арифметическое:
\[
\frac{18+54}{2}=\frac{72}{2}=36.
\]
Оба условия задачи выполнены.

\textcolor{teal!60!black}{Итог:} правильный ответ --- $18$ и $54$. Ошибка возникла именно при сложении $x+3x$: здесь должно быть $4x$.

\textbf{Задача на закрепление.}

Среднее арифметическое двух положительных чисел равно $40$. Одно число в $4$ раза больше другого. Найдите оба числа.

\newpage
\phantomsection\label{task:8}
\section*{Задание 8}

\textbf{Текст задания.}

Сумма трёх последовательных натуральных чисел равна $114$. Найдите эти числа.

\textbf{Ваше решение.}

В работе записано:
\[
3a:3=114,
\]
затем
\[
3a=114\cdot3,
\qquad
3a=342,
\qquad
a=342:3=114.
\]
После этого указаны три числа:
\[
114,
\qquad
115,
\qquad
116.
\]

\textbf{Ваш ответ.}
\[
114,
\quad
115,
\quad
116.
\]

\textcolor{red}{Ошибка:} неверно составлено исходное уравнение по условию о \emph{сумме} трёх последовательных чисел.

\textbf{Где именно возникает ошибка.}

Корректный элемент идеи --- обозначить одно из чисел буквой $a$. Но первый записанный математический шаг уже неверен:
\[
3a:3=114.
\]

Эта запись фактически сокращается до
\[
a=114,
\]
то есть число $114$ принимается за одно из искомых чисел. Однако по условию $114$ --- это \emph{сумма трёх чисел}, а не значение каждого из них и не их среднее арифметическое.

\textbf{Почему этот шаг неверен.}

Последовательные натуральные числа отличаются друг от друга на $1$. Если первое число обозначить через $a$, то три числа имеют вид
\[
a,
\qquad
a+1,
\qquad
a+2.
\]
Поэтому их сумма равна
\[
a+(a+1)+(a+2).
\]
Именно эта сумма по условию равна $114$:
\[
a+(a+1)+(a+2)=114.
\]

В записи $3a:3=114$ сразу возникают две проблемы. Во-первых, выражение $3a$ соответствует сумме трёх \emph{одинаковых} чисел $a+a+a$, но искомые числа не одинаковые: второе на $1$ больше первого, а третье на $2$ больше первого. Во-вторых, деление на $3$ относится к нахождению среднего арифметического, тогда как в условии дана сумма. Поэтому делить сумму трёх чисел на $3$ на этапе составления уравнения не нужно.

Из-за неверного исходного уравнения все дальнейшие вычисления уже относятся не к условию задачи. Проверка это сразу показывает:
\[
114+115+116=345\neq114.
\]

\textcolor{blue!60!black}{Ключевая идея:} для трёх последовательных натуральных чисел нужно записывать $a$, $a+1$, $a+2$ и приравнивать \emph{их сумму} к числу, данному в условии.

\textbf{Исправление.}

Обозначим первое число через $a$. Тогда:
\[
a+(a+1)+(a+2)=114.
\]
Раскрываем скобки и приводим подобные слагаемые:
\[
3a+3=114.
\]
Вычитаем $3$ из обеих частей:
\[
3a=111.
\]
Делим обе части на $3$:
\[
a=37.
\]
Тогда следующие два последовательных числа:
\[
a+1=38,
\qquad
a+2=39.
\]

\textbf{Полное правильное решение.}

Пусть первое из трёх последовательных натуральных чисел равно $a$. Тогда второе равно $a+1$, а третье --- $a+2$.

По условию их сумма равна $114$, поэтому составляем уравнение:
\[
a+(a+1)+(a+2)=114.
\]
Снимаем скобки:
\[
a+a+1+a+2=114.
\]
Складываем подобные слагаемые и числа $1$ и $2$:
\[
3a+3=114.
\]
Переносим $3$ в правую часть, то есть вычитаем $3$ из обеих частей:
\[
3a=114-3=111.
\]
Делим обе части на $3$:
\[
a=37.
\]
Следовательно, три числа:
\[
37,
\qquad
38,
\qquad
39.
\]

\textbf{Проверка.}

Числа действительно последовательные:
\[
38=37+1,
\qquad
39=38+1.
\]
Их сумма:
\[
37+38+39=75+39=114.
\]
Условие выполнено.

\textcolor{teal!60!black}{Итог:} правильный ответ --- $37$, $38$, $39$. В исходной записи условие о сумме было ошибочно заменено делением на $3$.

\textbf{Задача на закрепление.}

Сумма трёх последовательных натуральных чисел равна $96$. Найдите эти числа.

\vfill
\begin{center}
\small Лёвин Артём Александрович эксклюзивно для Нади Клименко
\end{center}

% Краткое сообщение Наде:
% Надя, основная часть работы выполнена уверенно. В задании 4 нужно внимательнее
% складывать подобные слагаемые, а в задании 8 --- отличать условие о сумме от
% условия о среднем арифметическом. После исправления этих двух моментов ход
% решения становится последовательным и легко проверяется подстановкой.

\end{document}
